% docs/math/learned_collision.tex
% Presentation deck (2026-11-13): the learned coarse-cell collision operator of branch
% neural-collision -- the ES-FP + cubic drift step and its learned outputs, the feature basis,
% training through the solver, the linear/Lasso model and the MLP sweep -- at the state of
% the lcfix / lccv runs.  Build: see docs/math/README.md (shared style shape_slides_style).
\documentclass[10pt,aspectratio=169]{beamer}
\newif\ifcodedeck
\codedeckfalse
% docs/math/shape_slides_style.tex
% Shared preamble of the shape-sensitivity decks (shape_derivation.tex,
% shape_derivation_code.tex, shape_validation.tex).  Each deck sets \codedecktrue or
% \codedeckfalse before % docs/math/shape_slides_style.tex
% Shared preamble of the shape-sensitivity decks (shape_derivation.tex,
% shape_derivation_code.tex, shape_validation.tex).  Each deck sets \codedecktrue or
% \codedeckfalse before % docs/math/shape_slides_style.tex
% Shared preamble of the shape-sensitivity decks (shape_derivation.tex,
% shape_derivation_code.tex, shape_validation.tex).  Each deck sets \codedecktrue or
% \codedeckfalse before \input{shape_slides_style}.
\usetheme{default}
\usecolortheme{dove}
\usefonttheme[onlymath]{serif}
\setbeamertemplate{navigation symbols}{}
\setbeamertemplate{footline}[frame number]
\setbeamerfont{frametitle}{size=\large}
\usepackage{amsmath,amssymb}
\usepackage{booktabs,array}
\usepackage{graphicx}
\usepackage{tikz}
\usetikzlibrary{arrows.meta,positioning}

\newcommand{\E}{\mathbb{E}}
\newcommand{\R}{\mathbb{R}}
\newcommand{\ind}{\mathbf{1}}
\newcommand{\dd}{\mathrm{d}}
\newcommand{\Lam}{\Lambda}
\newcommand{\pth}{\mathrm{post}}
\newcommand{\vin}{v^{\mathrm{in}}}
\newcommand{\vout}{v^{\mathrm{out}}}
\newcommand{\Lstar}{\mathcal L^*}
\newcommand{\jump}[1]{[\![#1]\!]}
\newcommand{\code}[1]{\texttt{#1}}

\definecolor{codebg}{RGB}{232,238,248}
\definecolor{good}{RGB}{0,110,0}
\definecolor{bad}{RGB}{170,0,0}
\definecolor{warn}{RGB}{190,100,0}
\newcommand{\ok}{\textcolor{good}{\textbf{yes}}}
\newcommand{\nil}{\textcolor{gray}{0}}

% \codenote{...}: the code references of a derivation frame.  Shown at the foot of the
% frame in the documentation deck, dropped from the presentation deck, so the two decks
% match frame for frame.
\newcommand{\codenote}[1]{%
  \ifcodedeck\par\vfill
  {\scriptsize\setlength{\fboxsep}{3pt}%
   \colorbox{codebg}{\parbox{\dimexpr\linewidth-2\fboxsep\relax}{\textbf{Code.} #1}}}%
  \fi}
% \presonly{...}: a line of the presentation deck only (spoken intuition that the
% documentation deck replaces with its code note).
\newcommand{\presonly}[1]{\ifcodedeck\else#1\fi}
.
\usetheme{default}
\usecolortheme{dove}
\usefonttheme[onlymath]{serif}
\setbeamertemplate{navigation symbols}{}
\setbeamertemplate{footline}[frame number]
\setbeamerfont{frametitle}{size=\large}
\usepackage{amsmath,amssymb}
\usepackage{booktabs,array}
\usepackage{graphicx}
\usepackage{tikz}
\usetikzlibrary{arrows.meta,positioning}

\newcommand{\E}{\mathbb{E}}
\newcommand{\R}{\mathbb{R}}
\newcommand{\ind}{\mathbf{1}}
\newcommand{\dd}{\mathrm{d}}
\newcommand{\Lam}{\Lambda}
\newcommand{\pth}{\mathrm{post}}
\newcommand{\vin}{v^{\mathrm{in}}}
\newcommand{\vout}{v^{\mathrm{out}}}
\newcommand{\Lstar}{\mathcal L^*}
\newcommand{\jump}[1]{[\![#1]\!]}
\newcommand{\code}[1]{\texttt{#1}}

\definecolor{codebg}{RGB}{232,238,248}
\definecolor{good}{RGB}{0,110,0}
\definecolor{bad}{RGB}{170,0,0}
\definecolor{warn}{RGB}{190,100,0}
\newcommand{\ok}{\textcolor{good}{\textbf{yes}}}
\newcommand{\nil}{\textcolor{gray}{0}}

% \codenote{...}: the code references of a derivation frame.  Shown at the foot of the
% frame in the documentation deck, dropped from the presentation deck, so the two decks
% match frame for frame.
\newcommand{\codenote}[1]{%
  \ifcodedeck\par\vfill
  {\scriptsize\setlength{\fboxsep}{3pt}%
   \colorbox{codebg}{\parbox{\dimexpr\linewidth-2\fboxsep\relax}{\textbf{Code.} #1}}}%
  \fi}
% \presonly{...}: a line of the presentation deck only (spoken intuition that the
% documentation deck replaces with its code note).
\newcommand{\presonly}[1]{\ifcodedeck\else#1\fi}
.
\usetheme{default}
\usecolortheme{dove}
\usefonttheme[onlymath]{serif}
\setbeamertemplate{navigation symbols}{}
\setbeamertemplate{footline}[frame number]
\setbeamerfont{frametitle}{size=\large}
\usepackage{amsmath,amssymb}
\usepackage{booktabs,array}
\usepackage{graphicx}
\usepackage{tikz}
\usetikzlibrary{arrows.meta,positioning}

\newcommand{\E}{\mathbb{E}}
\newcommand{\R}{\mathbb{R}}
\newcommand{\ind}{\mathbf{1}}
\newcommand{\dd}{\mathrm{d}}
\newcommand{\Lam}{\Lambda}
\newcommand{\pth}{\mathrm{post}}
\newcommand{\vin}{v^{\mathrm{in}}}
\newcommand{\vout}{v^{\mathrm{out}}}
\newcommand{\Lstar}{\mathcal L^*}
\newcommand{\jump}[1]{[\![#1]\!]}
\newcommand{\code}[1]{\texttt{#1}}

\definecolor{codebg}{RGB}{232,238,248}
\definecolor{good}{RGB}{0,110,0}
\definecolor{bad}{RGB}{170,0,0}
\definecolor{warn}{RGB}{190,100,0}
\newcommand{\ok}{\textcolor{good}{\textbf{yes}}}
\newcommand{\nil}{\textcolor{gray}{0}}

% \codenote{...}: the code references of a derivation frame.  Shown at the foot of the
% frame in the documentation deck, dropped from the presentation deck, so the two decks
% match frame for frame.
\newcommand{\codenote}[1]{%
  \ifcodedeck\par\vfill
  {\scriptsize\setlength{\fboxsep}{3pt}%
   \colorbox{codebg}{\parbox{\dimexpr\linewidth-2\fboxsep\relax}{\textbf{Code.} #1}}}%
  \fi}
% \presonly{...}: a line of the presentation deck only (spoken intuition that the
% documentation deck replaces with its code note).
\newcommand{\presonly}[1]{\ifcodedeck\else#1\fi}


\title{A learned collision operator for coarse-cell DSMC}
\subtitle{ES-Fokker--Planck with learned rates and cubic drift, trained through the solver}
\date{2026-11-13 \quad (intermediate state)}

\begin{document}
\begin{frame}
\titlepage
\end{frame}

% ---------------------------------------------------------------------------------------
\begin{frame}{1. The operator: ES-FP + reconstruction + learned cubic drift}
\begin{columns}[T]
\begin{column}{0.55\textwidth}
\small
In each coarse collision cell, every step:
\begin{itemize}\setlength{\itemsep}{2pt}
\item reconstructed target $u(x_i)=u_c+G_c(x_i-\bar x_c)$, peculiar $c_i=v_i-u(x_i)$;
  $\Pi=\langle cc^T\rangle$, $\theta=\tfrac13\operatorname{tr}\Pi$
\item rates \quad $r_s=e^{\delta_s}\,p/\mu(T)$, \quad $r_q=e^{\delta_q}\,\mathrm{Pr}\,p/\mu(T)$,
  \quad $\mathrm{Pr}=\tfrac23$
\item $a=e^{-r_q\Delta t/3}$, \; $e_s=e^{-r_s\Delta t}$ (raised where needed: Mathiaud--Mieussens
  $\nu$-clamp)
\end{itemize}
\vspace{-4pt}
\begin{align*}
c_i' &= a\,c_i + \operatorname{chol}(\Sigma)\,\xi_i
 + (1-a^3)\,\beta\Big(\tfrac{|c_i|^2}{\theta}-5\Big)c_i,\\
\Sigma &= \theta(1-a^2)I+(e_s-a^2)\,\Pi_{\mathrm{dev}},\qquad \xi_i\sim\mathcal N(0,I)
\end{align*}
then a projection restores each cell's momentum and energy exactly.
Exactly integrated: stress relaxes at $r_s$, heat flux at $r_q$, for any $\Delta t$; the
cubic term adds $\approx27(1-a^3)\beta\,q$ to the heat flux, $\approx0$ to stress.
\end{column}
\begin{column}{0.43\textwidth}
\footnotesize
\textbf{Learned per cell} (bounded, zero $\Rightarrow$ the physical operator):
\begin{center}
\begin{tabular}{@{}lll@{}}
\toprule
output & acts on & physics \\ \midrule
$\delta_s\in(-2,2)$ & $\log r_s$ & viscosity \\
$\delta_q\in(-2,2)$ & $\log r_q$ & conductivity \\
$\beta\in(-\tfrac12,\tfrac12)$ & cubic drift & heat flux via $\langle|c|^4c\rangle$ \\
\bottomrule
\end{tabular}
\end{center}
$(\delta_s,\delta_q,\beta)=B\tanh\!\big(f_\theta(\phi_c/\sigma)/B\big)$, $\phi_c$: 9 cell features.

\medskip
\textbf{Why:} coarse cells relax every particle toward one cell state
$\Rightarrow$ excess transport $\propto(\Delta x/\lambda)^2$
(NTC $+23\%$, ES-FP+recon $+11\%$ viscosity at $\Delta x=2\lambda$).
\textbf{Potential:} $\Delta x=\lambda/4\to2\lambda$ is $8\times$ per axis $=64\times$ fewer
particles in 2-D ($512\times$ in 3-D) at fixed particles/cell; a pathwise-differentiable,
conservative, Galilean/rotation-invariant operator.
\end{column}
\end{columns}
\end{frame}

% ---------------------------------------------------------------------------------------
\begin{frame}{2. The features: one term per known error mechanism}
\small
All local, dimensionless, invariant ($\lambda,\tau$ the cell's own VSS mean free path and time);
each standardised by its RMS over the training cells.
\begin{center}
\footnotesize
\renewcommand{\arraystretch}{1.12}
\begin{tabular}{@{}llp{0.6\textwidth}@{}}
\toprule
feature & form & mechanism \\ \midrule
dx2 & $(\Delta x/\lambda)^2$ & cell-size error of uniform in-cell collisions,
  $\mu_{\rm eff}/\mu-1=\tfrac{16}{45\pi}(\Delta x/\lambda)^2$ (Alexander--Garcia--Alder) \\
dx4 & $(\Delta x/\lambda)^4$ & next order of the same expansion \\
dt2 & $(\Delta t/\tau)^2$ & time-step error \\
dt1 & $\Delta t/\tau$ & residual first-order splitting error \\
inv\_n & $1/N_c$ & finite-particle bias of the cell state \\
dx2\_inv\_n & $(\Delta x/\lambda)^2/N_c$ & cell-size error measured with a noisy cell state \\
kn\_u2 & $(\lambda|S|/v_{\rm th})^2$ & Burnett-order (Chapman--Enskog $O(\mathrm{Kn}^2)$) stress
  closure \\
kn\_t2 & $(\lambda|\nabla\ln T|)^2$ & Burnett-order heat-flux closure \\
dx2\_kn2 & $(\Delta x/\lambda)^2(\mathrm{Kn}_u^2+\mathrm{Kn}_T^2)$ & cell-size error growing
  with gradients (the reconstruction is linear) \\
\bottomrule
\end{tabular}
\end{center}
\vspace{2pt}
\footnotesize
Not used: the 1-D collective (particles-per-$\lambda$) term -- non-local, suppressed in 2-D slabs;
a positive-definiteness-limiter term -- tried and dropped (it slows stress relaxation only
$\sim$2\% in the M10 core, and is mostly sampling noise).
\end{frame}

% ---------------------------------------------------------------------------------------
\begin{frame}{3. Learning through the solver}
\begin{columns}[T]
\begin{column}{0.5\textwidth}
\small
\textbf{Data} (no gradient): fine NTC at $\Delta x=\lambda/4$, 20 particles/cell --
4 shear waves ($u_0/v_{\rm th}\in\{0.1,0.5\}$, $L\in\{32,128\}\lambda$) and 4 normal shocks
(M 1.5, 3, 6, 10); coarse cases at $\Delta x=2\lambda$ (waves), $\lambda_1$ and $4\lambda_1$
(shocks). Train S3 S4 K1 K4 K8, test the other 7.

\medskip
\textbf{Window}: two independent coarse replicas, 5 blocks $\times$ 20 steps;
waves from thinned fine snapshots, shocks from \emph{persistent chains} at the model's own
steady state (snapshot starts hid the heat-flux deficit: $\int q_x$ 1.01 in the window vs 0.94
steady).

\medskip
\textbf{Deviation}: $d_{tkf}=P\big[F(m_t)-F(m^{\rm ref}_t)\big]_{kf}$, \; $F$: $P_{xy}$ (waves,
Fourier mode), $\tau_{xx},q_x$ per cell (shocks, reference placed at each replica's shock);
$w_f=1/\langle F_{\rm ref}^2\rangle$.
\end{column}
\begin{column}{0.48\textwidth}
\small
\textbf{Objective} $\tfrac12 w\,\E[d]^2$; running-bias estimator ($b\approx\E[d]$, stop-gradient):
\[
\nabla_\theta J=\operatorname*{mean}_{t,k}\sum_f w_f\, b_{tkf}\,\tfrac12
\Big(\tfrac{\partial m^a}{\partial\theta}+\tfrac{\partial m^b}{\partial\theta}\Big)_{tkf}
\]
\vspace{-4pt}
\begin{itemize}\setlength{\itemsep}{2pt}
\item $\partial m/\partial\theta$: \textbf{exact reverse-mode AD} through the 100 steps
  (checkpointed scan, keys replayed) -- collision (rates, Cholesky, cubic, projection,
  features), free flight, inflow; discrete events carry no gradient
\item cost $\approx$ a few forward windows; no fine-grid adjoint, no ensemble terms
\item cases summed, Adam with fan-in-normalised steps, tail-averaged iterates
\end{itemize}
\footnotesize
Checks: AD $=$ FD per window (common random numbers); estimator sign vs.\ steady FD at the base.
Open: the window start is held fixed (truncation of the steady-state response).
\end{column}
\end{columns}
\end{frame}

% ---------------------------------------------------------------------------------------
\begin{frame}{4. The linear model and the Lasso: which mechanisms matter}
\begin{columns}[T]
\begin{column}{0.52\textwidth}
\small
\textbf{Model}: $(\delta_s,\delta_q,\beta)=A\,\phi_c/\sigma$, \; $A\in\R^{3\times9}$
(18 coefficients without the cubic row). Each coefficient = the weight of one error mechanism
on one transport coefficient: directly interpretable.

\medskip
\textbf{Gauss--Newton Lasso}: per window $d(A)\approx d(A_0)+J(A-A_0)$, $J$ by forward-mode AD
(one tangent per coefficient), averaged per case:
\[
\min_A\ \tfrac12\sum_{\rm cases}\operatorname*{mean}_{\rm rows}
w\big(\bar d+\bar J(A-A_0)\big)^2+\alpha\|A\|_1
\]
coordinate descent along the $\alpha$ path; $\alpha$ by leave-one-case-out CV (one-s.e.\ rule);
stability selection over 200 window bootstraps; optional sign constraints
(cell-size/time-step terms may only speed up stress relaxation).
\end{column}
\begin{column}{0.46\textwidth}
\small
\textbf{First result} (snapshot windows, before chains):
\begin{itemize}\setlength{\itemsep}{2pt}
\item kept 2 of 18, stability 0.99--1.0: $1/N_c$ and $(\Delta x/\lambda)^2\mathrm{Kn}^2$, both
  speeding up relaxation (lower viscosity/conductivity); no Prandtl-row term
\item free and sign-constrained fits identical
\item squared bias $4.5\!\times\!10^{-3}\to3.0\!\times\!10^{-3}$
\item held-out (LOCO) error improves for S3, S4, K4, K8 at moderate $\alpha$, not K1;
  small $\alpha$ overfits (held-out K4 error $\times15$)
\end{itemize}
\footnotesize
To redo with chains and the cubic row; the linear fits also serve as the MLP's interpretable
baseline (linear $\approx$ MLP $\Rightarrow$ the theory basis suffices).
\end{column}
\end{columns}
\end{frame}

% ---------------------------------------------------------------------------------------
\begin{frame}{5. The MLP and its cross-validated design}
\begin{columns}[T]
\begin{column}{0.5\textwidth}
\small
\textbf{Model}: $f_\theta(\phi)=\mathrm{MLP}(\phi/\sigma)-\mathrm{MLP}(0)$, tanh hidden layers,
zero-initialised last layer $\Rightarrow$ starts as the physical operator and vanishes in
resolved equilibrium cells; same bounded outputs. Depth 2, width 32 $\approx$ 1.4k parameters.

\medskip
\textbf{Design sweep} (L9 orthogonal array, every pair of levels once):
\begin{center}
\begin{tabular}{@{}ll@{}}
\toprule
factor & levels \\ \midrule
hidden layers & 1, 2, 3 \\
width & 8, 32, 128 \\
weight decay & 0, $10^{-3}$, $10^{-2}$ \\
learning rate & $10^{-3}$, $3\!\times\!10^{-3}$, $10^{-2}$ \\
\bottomrule
\end{tabular}
\end{center}
\end{column}
\begin{column}{0.48\textwidth}
\small
\textbf{Per design point, 7 runs}: 5 leave-one-case-out folds (train on 4 of S3 S4 K1 K4 K8,
validate forward on the held-out one) and 2 all-case runs with different seeds (training
variance; seeds also pick each run's shock chains from the shared warm pool).
63 runs, $\sim$215 iterations each, sized to 55 GPU-h on 7 MIG slices (one shared warm pool,
built once).

\medskip
\textbf{Selection}: held-out error of the folds; then the chosen design trained on all five
and evaluated on the 7 test cases and a coarse M10 cylinder.

\medskip
\footnotesize
Steps: Adam divided by each layer's fan-in (lr = output change per iteration, independent of
width); decoupled weight decay; mean of the last 25\% of iterates.
\end{column}
\end{columns}
\end{frame}

% ---------------------------------------------------------------------------------------
\begin{frame}{Status (2026-11-13)}
\small
\begin{itemize}\setlength{\itemsep}{3pt}
\item \textbf{Base operators} (coarse, untrained): ES-FP+recon removes about half of NTC's
  cell-size viscosity error; its shock stress excess grows with Mach
  ($\int\tau_{xx}$ 1.05 at M1.5 $\to$ 1.17 at M10), the heat flux is 3--13\% low.
\item \textbf{Robust so far}: shear-wave viscosity corrected in every run
  (S3 1.065 $\to$ 0.97--1.05, S4 1.131 $\to$ 1.04--1.08); M10 stress partly (1.17 $\to$ 1.04--1.16).
\item \textbf{Diagnosed}: (i) snapshot-started windows never saw the steady heat-flux deficit
  $\Rightarrow$ persistent chains; (ii) chain pairs compared against one target at their mean
  shock position $\Rightarrow$ anti-correlated replicas (negative pair loss), a smeared-shock
  fake bias, and K4 worsening under its own gradient -- fixed (per-replica targets,
  bias table in the shock frame).
\item \textbf{Running}: re-checks with the fix (K4/K8 alone, linear/MLP, two rates vs.\ cubic),
  then the MLP sweep.
\item \textbf{Open}: window truncation (the start state is held fixed); what two rates can express
  in non-Gaussian shock cores (the cubic row); the weak M1.5 signal.
\end{itemize}
\end{frame}

\end{document}
